%\chapter*{\textit{\textbf{ABSTRACT}}}
\addcontentsline{toc}{chapter}{\textbf{ABSTRACT}}
%\chapter[1 cm]{title}
%\begin{center}
	%\textbf{{\large{ABSTRAK}}}
%\end{center}
%\cchapter{\textbf{\large{ABSTRAK}}}
%\begin{center}
%\textbf{{\LARGE{\textsc{MODEL PREDIKSI INAR(1) \\UNTUK STATUS KESEHATAN}}}}
%\end{center}
%\vspace{.8cm}
%\begin{center}
%\addcontentsline{toc}{chapter}{ABSTRAK}
%\end{center}

\begin{center}
\textbf{\fontsize{12pt}{0}\selectfont{\textit{Your Final Project Title}}} \\
\vspace{14pt}
\textbf{\fontsize{12pt}{0}\selectfont{Nama Mahasiswa (1xx45xxxx)}}\\
\textbf{\fontsize{12pt}{0}\selectfont{Nama Pembimbing 1 beserta gelar}}\\
\textbf{\fontsize{12pt}{0}\selectfont{Nama Pembimbing 2 beserta gelar}}\\
\vspace{36pt}
\textbf{\fontsize{12pt}{0}\selectfont{ABSTRACT}}
\end{center}
\vspace{20pt}
\singlespacing
The intended abstract is a summary or extract, a maximum of 200 words or one page. Abstract is an overview of a final project that contains problems, objectives, research methods, results, and conclusions. Abstracts are made to make it easier for readers to quickly understand the contents of the final project to decide whether to read further or not. Abstract does not contain pictures or tables, written in Times New Roman letters, 12 pt, one space. Abstracts are written in Indonesian and English, each starting on a new page. Abstract should contain one sentence explaining the background of the problem, one sentence explaining the objectives and one sentence explaining the benefits, one sentence explaining the scope and one sentence explaining the limitations of the problem, one sentence explaining the methodology, experiment and one sentence explaining the interpretation of the data as well as a sentence describing the research results obtained. Below the abstract, after a blank line, write "Keywords:" followed by the five appropriate keywords.

\singlespacing
\noindent 
\vspace{1ex}
\noindent
\begin{tabularx}{\textwidth}{@{}lX@{}}
    {\textbf{Keywords:}} & {Keyword 1, Keyword 2, Maximum 5 Keywords}
\end{tabularx}
